\documentclass[10pt,a4paper]{amsart}
\usepackage[margin=19mm]{geometry}
\usepackage{amsmath,amssymb,mathtools}
\usepackage[scr=rsfs,cal=euler]{mathalfa}
\usepackage{xfrac}
\usepackage[unicode,hidelinks,bookmarks=false]{hyperref}
\newcommand{\ie}{i.e.\ }
\newcommand{\cf}{cf.\ }
\newcommand{\resp}{respectively\ }
\newcommand{\Subsection}{\subsection}
\providecommand{\nobreakdash}{\nobreak\hbox{-}}
\setlength{\emergencystretch}{2em}
\multlinegap0pt
\usepackage{amssymb}
\newtheorem{theorem}{Theorem}
\title{Uniqueness in Lorentz Spaces of the 2d Navier-Stokes Equations}
\author{Alexandru F. Radu}
\date{Source: arXiv:2603.02354v1 (CC BY 4.0)}
\begin{document}
\maketitle

\noindent\emph{Source and licence.} Uniqueness in Lorentz Spaces of the 2d Navier-Stokes Equations. Version arXiv:2603.02354v1 (CC BY 4.0), distributed by arXiv under the Creative Commons Attribution 4.0 International licence, http://creativecommons.org/licenses/by/4.0/, as recorded in the arXiv licence metadata for this exact version and shown on the abstract page https://arxiv.org/abs/2603.02354v1. Full licence notice: \texttt{licenses/arxiv-2603.02354-CC-BY-4.0.txt}.\par\smallskip

\noindent\textit{Editorial scope.} Counted unit: the article's theorem main (source0.tex lines 426--442). The article's complete original proof of it is delivered with it (source0.tex lines 444--563, one \texttt{proof} environment of 120 lines). No mathematics is written, repaired or re-derived: the statement and the proof are the article's own lines, in the article's own notation, and no retained line is substituted or re-flowed.  No further source-local context is needed: every cross-reference of every retained line resolves inside the two delivered spans.  Nothing was omitted from any retained span, so the package carries no omission record.  No retained line contains a citation, so the wrapper carries no bibliography and no bibliography entry.  No retained line contains a figure environment, an image-inclusion command or drawing code, so no image is delivered and the package declares no asset dependency.  Editorial formatting only, in addition: an AMS \texttt{amsart} preamble that restates the article's own declarations and macros used by the retained lines, namely the environments \texttt{theorem} on separate counters, together with the standard packages those lines need.  The retained environments print in this document as Theorem 1 in order of appearance, whereas the article prints the same items with the numbering of its own full document.  The complete native source of the article as distributed in the arXiv e-print for this exact version is bundled unchanged as \texttt{sources/195613/source0.tex}.
\begin{theorem}\label{main}
Let $T\in(0,\infty]$ and let $1<q<2$. Let $v_1,v_2$ be mild solutions of the periodic Navier--Stokes equation on
$[0,T)\times\mathbb T^2$ with the same initial data
\[
v_1(0)=v_2(0)\in L^{2,q}(\mathbb T^2).
\]
Assume:
\begin{enumerate}
\item $v_j\in C([0,T),L^{2,q}(\mathbb T^2))$ for $j=1,2$.
\item (Short-time smoothing at every restart time) For every $T_0\in[0,T)$,
\begin{equation}\label{eq:S}
\lim_{\delta\downarrow 0}\ \sup_{t\in(T_0,T_0+\delta]}
\sqrt{t-T_0}\,\|v_j(t)\|_{L^\infty(\mathbb T^2)}=0,\qquad j=1,2.
\end{equation}
\end{enumerate}
Then $v_1\equiv v_2$ on $[0,T)$.
\end{theorem}

\begin{proof}
\textbf{Step 1.}
Since $1<q<2$, the Lorentz embedding monotonicity in the second index yields
\[
L^{2,q}(\mathbb T^2)\hookrightarrow L^{2,2}(\mathbb T^2)=L^2(\mathbb T^2)
\]
continuously. Hence $v_j\in C([0,T),L^{2,q})$ implies $v_j\in C([0,T),L^2)$, and therefore
\[
w:=v_1-v_2\in C([0,T),L^2(\mathbb T^2)).
\]

\textbf{Step 2.}
Define
\[
T^*:=\sup\{S\in[0,T): w(t)=0\ \text{in }L^2 \text{ for all }t\in[0,S)\}.
\]
If $T^*=T$ we are done. Assume $T^*<T$. Since $w\in C([0,T),L^2)$ and $w(t)=0$ for $t<T^*$, we have
\begin{equation}\label{eq:wTstar}
w(T^*)=0\quad\text{in }L^2(\mathbb T^2).
\end{equation}

\textbf{Step 3.}
By the time-restart property of mild solutions, for each $j=1,2$ and all $t\in[T^*,T)$,
\begin{equation}\label{eq:mild-restart}
v_j(t)=e^{(t-T^*)\Delta}v_j(T^*)
-\int_{T^*}^{t} K_{\mathrm{per}}(t-s)*\bigl(v_j(s)\otimes v_j(s)\bigr)\,ds,
\end{equation}
where $K_{\mathrm{per}}$ is the periodic Oseen-type kernel associated to the Leray projector. Subtracting the identities \eqref{eq:mild-restart} for $j=1$ and $j=2$,
the linear terms cancel because of \eqref{eq:wTstar}, hence
\begin{equation}\label{eq:w-duhamel}
w(t)=-\int_{T^*}^{t} K_{\mathrm{per}}(t-s)*
\Bigl(v_1(s)\otimes v_1(s)-v_2(s)\otimes v_2(s)\Bigr)\,ds.
\end{equation}
Using the algebraic identity
\[
v_1\otimes v_1-v_2\otimes v_2=(v_1-v_2)\otimes v_1+v_2\otimes(v_1-v_2)=w\otimes v_1+v_2\otimes w,
\]
we rewrite \eqref{eq:w-duhamel} as
\begin{equation}\label{eq:w-duhamel2}
w(t)=-\int_{T^*}^{t} K_{\mathrm{per}}(t-s)*
\Bigl(w(s)\otimes v_1(s)+v_2(s)\otimes w(s)\Bigr)\,ds.
\end{equation}

\textbf{Step 4.}
We use the kernel estimate
\begin{equation}\label{eq:K-L1}
\|K_{\mathrm{per}}(t)\|_{L^1(\mathbb T^2)}\le C\,t^{-1/2},\qquad t>0.
\end{equation}
Fix $t\in[T^*,T)$. Taking $L^2$ norms in \eqref{eq:w-duhamel2} and using the triangle inequality gives
\begin{align}
\|w(t)\|_{2}
&\le \int_{T^*}^{t}\bigl\|K_{\mathrm{per}}(t-s)*\bigl(w(s)\otimes v_1(s)\bigr)\bigr\|_{2}\,ds
\nonumber\\
&\quad +\int_{T^*}^{t}\bigl\|K_{\mathrm{per}}(t-s)*\bigl(v_2(s)\otimes w(s)\bigr)\bigr\|_{2}\,ds.
\label{eq:basic-tri}
\end{align}
By Young's inequality $L^1*L^2\to L^2$ and \eqref{eq:K-L1},
\[
\|K_{\mathrm{per}}(t-s)*F\|_2\le \|K_{\mathrm{per}}(t-s)\|_1\|F\|_2\le C(t-s)^{-1/2}\|F\|_2,
\]
hence from \eqref{eq:basic-tri}
\begin{align}
\|w(t)\|_2
&\le C\int_{T^*}^{t}(t-s)^{-1/2}\,\|w(s)\otimes v_1(s)\|_2\,ds
+ C\int_{T^*}^{t}(t-s)^{-1/2}\,\|v_2(s)\otimes w(s)\|_2\,ds.
\label{eq:after-young}
\end{align}
Using the product estimate $L^\infty\times L^2\subset L^2$,
\[
\|w(s)\otimes v_1(s)\|_2\le \|v_1(s)\|_\infty\|w(s)\|_2,\qquad
\|v_2(s)\otimes w(s)\|_2\le \|v_2(s)\|_\infty\|w(s)\|_2,
\]
we obtain
\begin{equation}\label{eq:Volterra}
\|w(t)\|_2\le C\int_{T^*}^{t}(t-s)^{-1/2}\bigl(\|v_1(s)\|_\infty+\|v_2(s)\|_\infty\bigr)\|w(s)\|_2\,ds.
\end{equation}

\textbf{Step 5.}
Fix $\delta\in(0,\min\{1,T-T^*\})$ and restrict to $t\in[T^*,T^*+\delta]$. Define
\begin{equation}\label{eq:Mdelta}
M(\delta):=\sup_{s\in(T^*,T^*+\delta]}\sqrt{s-T^*}\,\bigl(\|v_1(s)\|_\infty+\|v_2(s)\|_\infty\bigr).
\end{equation}
Then for $s\in(T^*,T^*+\delta]$,
\begin{equation}\label{eq:bound-vinf}
\|v_1(s)\|_\infty+\|v_2(s)\|_\infty\le \frac{M(\delta)}{\sqrt{s-T^*}}.
\end{equation}
By the assumption \eqref{eq:S} applied at $T_0=T^*$ for each $v_j$, we have
\begin{equation}\label{eq:Mto0}
\lim_{\delta\downarrow 0} M(\delta)=0.
\end{equation}
Next, set
\[
\|w\|_{L^2_\delta}:=\sup_{t\in[T^*,T^*+\delta]}\|w(t)\|_2.
\]
Using \eqref{eq:Volterra}, \eqref{eq:bound-vinf}, and the definition of $\|w\|_{L^2_\delta}$ yields, for $t\in[T^*,T^*+\delta]$,
\begin{equation}\label{eq:w-est-preBeta}
\|w(t)\|_2
\le C\,M(\delta)\,\|w\|_{L^2_\delta}\int_{T^*}^{t}(t-s)^{-1/2}(s-T^*)^{-1/2}\,ds.
\end{equation}
Compute the integral explicitly. Let $a:=t-T^*>0$ and set $s=T^*+a\theta$, $ds=a\,d\theta$, $\theta\in[0,1]$:
\begin{align*}
\int_{T^*}^{t}(t-s)^{-1/2}(s-T^*)^{-1/2}\,ds
&=\int_0^1 (a(1-\theta))^{-1/2}(a\theta)^{-1/2}\,a\,d\theta\\
&=\int_0^1 \theta^{-1/2}(1-\theta)^{-1/2}\,d\theta\\
&=\mathrm{B}\Bigl(\frac12,\frac12\Bigr)=\pi.
\end{align*}
Hence \eqref{eq:w-est-preBeta} becomes
\[
\|w(t)\|_2\le C\pi\,M(\delta)\,\|w\|_{L^2_\delta}\qquad\forall t\in[T^*,T^*+\delta].
\]
Taking the supremum over $t\in[T^*,T^*+\delta]$ yields
\begin{equation}\label{eq:contraction}
\|w\|_{L^2_\delta}\le C\pi\,M(\delta)\,\|w\|_{L^2_\delta}.
\end{equation}

\textbf{Step 6.}
By \eqref{eq:Mto0}, choose $\delta>0$ so small that $C\pi\,M(\delta)<1$. Then \eqref{eq:contraction} implies
$\|w\|_{L^2_\delta}=0$, i.e.\ $w(t)=0$ in $L^2$ for all $t\in[T^*,T^*+\delta]$.
This contradicts the definition of $T^*$ unless $T^*=T$. Therefore $T^*=T$ and $v_1\equiv v_2$ on $[0,T)$.
\end{proof}

\end{document}
